% page 457 of the facsimile (image p-456 of the scan)
% block 165.1 x 237.7 mm ; left margin 36.9 mm
\pgc{15.240mm}{0.000mm}{
\l{\cel{54}449}
\l{}
\l{\cel{1}\sou{oodoido}. (\sou{matem.}) Singla de la kurvi quin onu obtenas per}
\l{\cel{4}prenar la simetrii di un punto Q, relate la tangenti di}
\l{\cel{4}ta kurvo.}
\l{}
\l{\cel{1}\sou{poemo}. Versaro pasable longa. - DEFIRS.}
\l{}
\l{poeto. Persono qua su konsakras a poezio ; persono qua esass}
\l{\cel{4}dotita por poezio. - DEFIRS.}
\l{}
\l{\cel{1}\sou{poetiko}. Verko docala, en qua expozesas la reguli di la kompozo}
\l{\cel{4}per versi. - DEFIS.}
\l{}
\l{\cel{1}\sou{poezio}. I. La arto facar verki ek versi. - II.BHeleso di la ma-}
\l{\cel{4}terio e di la expreso-maniero, maxim grava pri poezio, sen-}
\l{\cel{4}depende de la versifo. - DEFIRS.}
\l{}
\l{\cel{1}\sou{poka}. Mikra-quanta. - FIS.}
\l{}
\l{\cel{1}\sou{polo}.- I. (\sou{astron}.) Singla ek la extremaji di la axo di mondo,}
\l{\cel{4}cirkum qua semblas movar su, en 24 hori, la sfero ciela. -}
\l{\cel{4}II. (\sou{geom.}) Singla ek la extremaji di la axo cirkum qua,}
\l{\cel{4}supozite, rotacas korpo sfera, elipsatra. - III. (\sou{geogr}.)}
\l{\cel{4}Singla ek la du extremaji di la axo di Tero, cirkum qua la}
\l{\cel{4}Terglobo semblas rotacar en 24 hori. - IV. (\sou{fiziko)}. (a)}
\l{\cel{4}Singla ek la du punti inter-opozita di magneto ube eventas}
\l{\cel{3}maxima-force la atrakti e la repulsi magnetala, (b) Singlaa}
\l{\cel{3}ek la du extremaji inter-opozita di un fluo elektrala. --}
\l{\cel{4}DEFIRS.}
\l{}
\l{\sou{polarimetro}. (\sou{fiziko)}. Instrumento por mezurar la proporciono}
\l{\cel{4}di la lumo palarigita qua existas en lumo-radio o la rotacodi}
\l{\cel{4}la polarigo-plano. -}
\l{}
\l{\cel{1}\sou{polemikar}. (\sou{netrans., kun ulu, kontre ulu}). Debatar skribe.}
\l{\cel{4}- DEFIRS.}
\l{}
\l{\cel{1}\sou{poleno}. (\sou{bot.}) Polvo fekundigiva che la vejetanti. - DEFIS.}
\l{}
\l{\cel{1}\sou{polexo}. (\sou{anat.})\cel{2}La fingro maxim grosa e maxim forta di la}
\l{\cel{3}manui e di la pedi. - EIL..}
\l{}
\l{\cel{1}\sou{poli..}. Sufixo ciencala : \textquotedbl{}multa-....a.}
\l{}
\l{\cel{1}\sou{poliandra}. Spozino qua havas plura spozuli. - DEFIRS.}
\l{}
\l{\cel{1}\sou{polico.}\cel{2}I. Regulari facita por mantenar la ordino e la seku-}
\l{\cel{4}reso statala, nacionala. - II. Surveyo quan onu agas por igr}
\l{\cel{4}aplikata la regulari, por la manteno di la ordino e di la}
\l{\cel{4}sekureso statala, nacionala. - DEFIRS.}
\l{}
\l{\cel{1}\sou{polidaktila.}\cel{1}(\sou{teratol.)}\cel{1}Qua havas fingri supernombra. -}
\l{}
\l{\cel{1}poliedro. (\sou{geom.)}\cel{1}Korpo solida qua havas plura edri.DEFIRS.}
}
